\documentclass[12pt,a4paper]{article}

% --- Paquetes Básicos ---
\usepackage[utf8]{inputenc}
\usepackage[spanish,es-tabla]{babel}
\usepackage[margin=2.5cm]{geometry}
\usepackage{graphicx}
\usepackage{titlesec}
\usepackage{enumitem}
\usepackage{hyperref}

% --- Configuración de Títulos ---
\titleformat{\section}{\normalfont\Large\bfseries}{\thesection}{1em}{}
\titleformat{\subsection}{\normalfont\large\bfseries}{\thesubsection}{1em}{}

\begin{document}

% --- Portada ---
\begin{titlepage}
    \centering
    \includegraphics[width=0.4\textwidth]{logo_unan.png}\\[1cm] % Reemplazar con tu imagen
    {\scshape\LARGE Universidad Nacional Autónoma de Nicaragua \par}
    {\scshape\Large CUR-Carazo \par}
    \vspace{1cm}
    {\large Facultad de Ingeniería Tecnología y Salud \par}
    {\large Ingeniería en Sistemas \par}
    \vspace{1.5cm}
    {\huge\bfseries Propuesta de diseño conceptual de un sistema de gestión de turnos en centros de salud \par}
    \vspace{0.5cm}
    {\large Análisis de la problemática, modelado conceptual y evaluación del prototipo \par}
    \vspace{1.5cm}
    
    \textbf{Trabajo académico de las asignaturas:} \\
    Integrador I, Diseño de Interfaces de Usuario e \\ Introducción a la programación \par
    \vspace{1cm}
    
    \textbf{Autores:} \\
    Emily Nazareth Murillo Montenegro \\
    Cristofer Francisco Blass Espinoza \\
    Abel Ezequiel Cerda Martínez \\
    Pedro Luis Rodríguez López \par
    \vspace{1cm}
    
    \textbf{Tutor:} \\
    Msc. Martha Rodríguez \par
    \vfill
    {\large Carazo, Nicaragua \\ 2026 \par}
\end{titlepage}

\tableofcontents
\newpage

% --- Introducción ---
\section{Introducción}
En la actualidad, muchos centros de salud de San Marcos Carazo enfrentan dificultades en la organización de la atención[cite: 30]. El manejo manual de turnos genera largas filas y falta de priorización[cite: 30, 31]. Se propone un Sistema de Gestión de Turnos para automatizar el registro y la asignación según niveles de prioridad[cite: 32, 33].

% --- Justificación ---
\section{Justificación}
El sistema surge para mejorar la eficiencia médica[cite: 35]. El proceso manual actual genera estrés laboral y disminuye la calidad del servicio[cite: 39]. La automatización permitirá una atención justa para emergencias, adultos mayores y mujeres embarazadas[cite: 40, 41].

% --- Objetivos ---
\section{Objetivos}
\subsection{Objetivo General}
Desarrollar un prototipo funcional de un sistema de gestión de turnos con asignaciones automáticas para reducir tiempos de espera[cite: 43, 44, 46].

\subsection{Objetivos Específicos}
\begin{itemize}
    \item Crear un módulo para que el personal de salud gestione la lista de pacientes[cite: 49].
    \item Diseñar un registro de pacientes con captura de datos y condición médica[cite: 50].
    \item Clasificar automáticamente a los pacientes según prioridad (emergencia, media, normal)[cite: 51].
    \item Implementar el cálculo de tiempo estimado de atención[cite: 52].
\end{itemize}

% --- Desarrollo del Proyecto ---
\section{Desarrollo del Proyecto}
\subsection{Conceptualización}
Una herramienta diseñada para administrar la atención mediante la optimización del tiempo y la equidad en el servicio[cite: 56, 57].

\subsection{Análisis del Problema}
La ausencia de sistemas automatizados provoca:
\begin{itemize}
    \item Desorden en la atención[cite: 64].
    \item Falta de priorización de casos urgentes[cite: 65].
    \item Sobrecarga operativa para el personal médico[cite: 67, 68].
\end{itemize}

% --- Diagnóstico y Árbol del Problema ---
\section{Diagnóstico del Problema}
\subsection{Árbol del Problema}
\textbf{Problema Central:} Deficiencias en la gestión de turnos[cite: 78].
\begin{itemize}
    \item \textbf{Causas:} Turnos manuales y demora en la atención[cite: 80, 82].
    \item \textbf{Consecuencias:} Pacientes molestos y errores en turnos[cite: 85, 86].
\end{itemize}

% --- Mapa de Empatía ---
\section{Mapa de Empatía}
\begin{description}
    \item[¿Qué dice?] "Llevo mucho tiempo esperando", "Esto está desorganizado"[cite: 93, 96].
    \item[¿Qué piensa?] Que el sistema es injusto y su tiempo no se respeta[cite: 98, 100].
    \item[¿Qué siente?] Frustración, impaciencia y estrés[cite: 113, 115].
\end{description}

% --- Planteamiento del Problema ---
\section{Planteamiento del Problema}
La dependencia de métodos manuales y papel ralentiza la recepción de usuarios y propicia errores humanos[cite: 128, 129]. Esto vulnera la equidad del servicio al no clasificar objetivamente las urgencias médicas[cite: 134].

% --- Anexos ---
\section{Anexos}
\begin{figure}[h]
    \centering
    % \includegraphics[width=0.8\textwidth]{diagrama_flujo.png}
    \caption{Diagrama de flujo del sistema de gestión de turnos.}
\end{figure}


% --- Arquitectura y Diseño del Sistema ---
\section{Arquitectura y Diseño del Sistema}
La solución se basa en una arquitectura de gestión de flujos de trabajo en tiempo real. El sistema procesa la entrada del paciente, aplica un algoritmo de priorización y gestiona una cola dinámica que el personal médico puede manipular[cite: 58, 135].

\subsection{Lógica de Priorización}
El sistema implementa un motor de reglas basado en las siguientes categorías de prioridad:
\begin{enumerate}
    \item \textbf{Prioridad 1 (Emergencia):} Atención inmediata, interrumpe el flujo normal de turnos.
    \item \textbf{Prioridad 2 (Alta Prioridad):} Atención preferente para adultos mayores y mujeres embarazadas.
    \item \textbf{Prioridad 3 (Normal):} Atención por orden de llegada para consultas generales.
\end{enumerate}

% --- Modelado de Datos ---
\section{Modelado de Datos}
Para el funcionamiento del prototipo, se han identificado las siguientes entidades clave dentro de la base de datos[cite: 50, 52]:

\begin{table}[h]
\centering
\begin{tabular}{|l|l|p{8cm}|}
\hline
\textbf{Entidad} & \textbf{Atributos Principales} & \textbf{Descripción} \\ \hline
Paciente & ID, Nombre, Condición Médica & Almacena los datos básicos y clínicos del usuario. \\ \hline
Turno & Número, Hora Entrada, Prioridad & Gestiona el puesto en la fila y el tiempo estimado[cite: 52]. \\ \hline
Personal & Usuario, Rol, Consultorio & Permite el acceso seguro al panel de gestión[cite: 49]. \\ \hline
\end{tabular}
\caption{Diccionario de datos básico del sistema.}
\end{table}

% --- Diseño de Interfaces de Usuario ---
\section{Diseño de Interfaces de Usuario}
El diseño se centra en la usabilidad tanto para el paciente como para el personal administrativo, buscando reducir la frustración identificada en el mapa de empatía[cite: 113, 127].

\subsection{Módulo del Paciente (Kiosco)}
Consiste en una interfaz táctil simplificada donde el usuario ingresa su nombre y selecciona su condición médica. El sistema genera un comprobante con el número de turno y el tiempo de espera calculado[cite: 50, 52].

\subsection{Panel del Personal de Salud}
Una interfaz de escritorio que permite visualizar la lista de espera ordenada automáticamente por prioridad. Incluye funciones para[cite: 49]:
\begin{itemize}
    \item Llamar al siguiente turno (activación del sistema de sonido y pantalla).
    \item Marcar paciente como atendido.
    \item Generar reportes de flujo y tiempos de espera.
\end{itemize}

% --- Evaluación del Prototipo ---
\section{Evaluación del Prototipo y Reportes}
El sistema genera datos estadísticos para la toma de decisiones, tales como la cantidad de pacientes atendidos, el tiempo promedio de espera por categoría y los tipos de consulta más frecuentes[cite: 42, 135].

\end{document}